% Book 3, Year 1, chapter 22 -- greedy algorithms and their proofs.
% Provenance (never printed): sources/book3/BRIEFS.md, chapter 22.
% ledger: b1:huffman-1952
\chapter{Greedy Algorithms and Their Proofs}\label{ch:b1:greedy-proofs}

A conference has forty talks to fit in one room, and each talk has fixed hours. Taking the shortest talks first sounds
reasonable; so does taking the earliest. One rule always gives the most talks, and it is neither of these. \omterm{def:b1:greedy-proofs:greedy}{Greedy
algorithms} build a solution one irrevocable choice at a time; they are fast and short, and they are often wrong. This
chapter is about telling the two cases apart, with two kinds of proof, and about one \omterm{def:b1:greedy-proofs:greedy}{greedy algorithm} that every
compressed file relies on.

\begin{recall}
A \omterm{def:b1:heaps:pq}{priority queue} (\cref{def:b1:heaps:pq}) gives its smallest element in $O(\log n)$. In a \omterm{def:b1:induction-and-termination:binary-tree}{binary tree}
(\cref{def:b1:induction-and-termination:binary-tree}), the depth of a node is the number of edges from the root to it.
\end{recall}

\section{The greedy method}\label{sec:b1:greedy-proofs:method}

\begin{definition}[Greedy algorithm]\label{def:b1:greedy-proofs:greedy}
A \emph{greedy algorithm}\index{greedy algorithm} builds a solution by a sequence of choices, each made by a fixed rule
that looks only at the current \omterm{def:b1:backtracking:backtracking}{partial solution}, and never undone.
\end{definition}

The second book met greedy methods on examples; here they need proofs. \omterm{def:b1:backtracking:backtracking}{Backtracking} (\cref{ch:b1:backtracking}) tries
every choice and undoes the bad ones; a \omterm{def:b1:greedy-proofs:greedy}{greedy algorithm} makes one choice per step and keeps it, which is why it is
fast, and why a proof is needed: nothing in the method guarantees that a choice good now is not fatal later.

\begin{method}[Proving a greedy algorithm optimal]\label{met:b1:greedy-proofs:exchange}
Two templates.
\begin{itemize}
\item \emph{Exchange argument}\index{exchange argument}: take any optimal solution; show that it can be changed, step by
step, into the greedy one without getting worse. The step usually replaces the optimal solution's first choice by the
greedy choice and checks that the result is still valid and no worse.
\item \emph{Greedy stays ahead}: measure the progress of a \omterm{def:b1:backtracking:backtracking}{partial solution}, and show by induction that after $i$
choices the greedy solution is at least as far as any other solution after $i$ choices.
\end{itemize}
Either proof must use the rule: a proof that would work for any rule proves nothing, since most rules fail.
\end{method}

\section{Interval scheduling}\label{sec:b1:greedy-proofs:intervals}

\begin{definition}[Interval scheduling, interval partitioning, maximum overlap]\label{def:b1:greedy-proofs:intervals}
Talks are intervals $[s, f)$ of time; two talks are \emph{compatible} if their intervals do not overlap. The
\emph{interval scheduling problem}\index{interval scheduling problem} asks for a largest set of pairwise compatible
talks. The \emph{interval partitioning problem}\index{interval partitioning} asks to place all the talks in as few rooms
as possible, compatible talks only in one room. The \emph{maximum overlap}\index{maximum overlap} of a set of talks is
the largest number of them going on at one moment.
\end{definition}

Four rules come to mind: earliest start, shortest, fewest conflicts with the others, and earliest finish. The program
applies each, keeping every talk compatible with those already kept, and compares with the optimum found by trying all
subsets.

\omcode{code/bachelor-1/22-greedy-proofs/listings/intervals.py}{6}{25}{Four greedy rules for one
room}\label{lst:b1:greedy-proofs:intervals}

\omoutput{code/bachelor-1/22-greedy-proofs/outputs/intervals.txt}{The conference, and the four rules against the
optimum on random sets of talks.}\label{out:b1:greedy-proofs:intervals}

A rule that fails once is wrong, however often it succeeds: fewest conflicts fails on one set in five hundred here.
Earliest start fails on one long talk that starts first; shortest fails on a short talk overlapping two long compatible
ones.

\begin{theorem}[Earliest finish is optimal]\label{thm:b1:greedy-proofs:earliest-finish}
Taking the talks by increasing finishing time, and keeping each one compatible with those kept, gives a largest set of
compatible talks.
\end{theorem}

\begin{proof}
By the \omterm{met:b1:greedy-proofs:exchange}{exchange argument}, by induction on the number of talks. Let $g$ be the talk that finishes first, which the greedy
keeps, and $O$ an optimal set, sorted by finishing time, with first talk $o$. Then $f(g) \leq f(o)$, and every other
talk of $O$ starts after $o$ finishes, hence after $g$ finishes: $O' = O - \{o\} + \{g\}$ is compatible, of the same
size, and optimal. The other talks of $O'$ all start at or after $f(g)$: they form a compatible set of the talks
starting at or after $f(g)$, and no larger one exists, or adding $g$ to it would beat $O$. After keeping $g$, the
greedy discards exactly the talks that start before $f(g)$ and runs on the others, where by induction it finds an
optimal set. With $g$, it has $|O|$ talks.
\end{proof}

\begin{omfigure}
\begin{tikzpicture}[x=12mm, y=7mm, t/.style={draw, rounded corners=1pt, minimum height=4mm, inner sep=0pt,
  font=\scriptsize}]
\draw[->] (0,-0.6) -- (10.4,-0.6) node[right, font=\scriptsize] {time};
\foreach \x in {0,2,...,10} \draw (\x,-0.5) -- (\x,-0.7) node[below, font=\scriptsize] {\x};
\node[font=\scriptsize, anchor=east] at (-0.1,3.6) {all talks};
\foreach \s/\f/\n/\r in {0/4/A/3.9, 1/3/B/3.3, 3/6/C/3.3, 4/7/D/3.9, 5/9/E/4.5, 7/8/F/3.9, 8/10/G/3.3}
  \node[t, fill=gray!15, minimum width={(\f-\s)*12mm}, anchor=west] at (\s,\r) {\n};
\node[font=\scriptsize, anchor=east] at (-0.1,2.2) {greedy};
\foreach \s/\f/\n in {1/3/B, 3/6/C, 7/8/F, 8/10/G}
  \node[t, fill=omDef!25, minimum width={(\f-\s)*12mm}, anchor=west] at (\s,2.2) {\n};
\node[font=\scriptsize, anchor=east] at (-0.1,1.2) {an optimum $O$};
\foreach \s/\f/\n in {0/4/A, 4/7/D, 7/8/F, 8/10/G}
  \node[t, fill=omThm!15, minimum width={(\f-\s)*12mm}, anchor=west] at (\s,1.2) {\n};
\node[font=\scriptsize, anchor=east] at (-0.1,0.2) {$O$ after one exchange};
\foreach \s/\f/\n in {1/3/B, 4/7/D, 7/8/F, 8/10/G}
  \node[t, fill=omThm!15, minimum width={(\f-\s)*12mm}, anchor=west] at (\s,0.2) {\n};
\draw[->, omThm, thick] (2,1.0) -- (2,0.45);
\end{tikzpicture}
\par\smallskip{\small Seven talks, the greedy choice by earliest finish, an optimal set, and the exchange step: the
optimum's first talk A is replaced by B, which finishes earlier, so the rest of the optimum still fits after it.}
\end{omfigure}

With several rooms, every talk must be placed. At a moment when $d$ talks are going on, $d$ rooms are needed, so the
\omterm{def:b1:greedy-proofs:intervals}{maximum overlap} is a lower bound; a greedy assignment reaches it.

\omcodes{code/bachelor-1/22-greedy-proofs/listings/partition_rooms.py}{9-16,19-30}{The \omterm{def:b1:greedy-proofs:intervals}{maximum overlap}, and rooms
assigned by increasing start}\label{lst:b1:greedy-proofs:rooms}

\omoutput[2-4]{code/bachelor-1/22-greedy-proofs/outputs/partition_rooms.txt}{The forty talks in nine rooms; the first
two rooms shown.}\label{out:b1:greedy-proofs:rooms}

\begin{proposition}[Rooms by increasing start]\label{prop:b1:greedy-proofs:rooms}
Assigning the talks by increasing start, each to a room whose last talk is over if there is one, uses exactly as many
rooms as the \omterm{def:b1:greedy-proofs:intervals}{maximum overlap}.
\end{proposition}

\begin{proof}
At least that many are needed. Conversely, suppose the algorithm opens room number $r + 1$ for a talk starting at $s$.
Then each of the $r$ rooms has a last talk that started at or before $s$ (talks come by increasing start) and is not
over at $s$: with the new talk, $r + 1$ talks are going on at time $s$, so $r + 1$ is at most the \omterm{def:b1:greedy-proofs:intervals}{maximum overlap}.
\end{proof}

\section{Unit tasks with deadlines and penalties}\label{sec:b1:greedy-proofs:penalties}

\begin{definition}[Unit-task scheduling problem]\label{def:b1:greedy-proofs:unit-tasks}
In the \emph{unit-task scheduling problem}\index{unit-task scheduling problem}, $n$ tasks each take one slot of time on
one machine, in slots $1, 2, \dots, n$; task $i$ has a deadline $d_i$ and a penalty $p_i > 0$, paid if it ends after slot
$d_i$. The problem asks for an order with the smallest total penalty.
\end{definition}

Only the set of tasks done on time matters: they can be run by increasing deadline, the late ones after them in any
order. A set is \emph{feasible} if its tasks can all be on time; minimising the penalty paid means maximising the
penalty of a feasible set.

\begin{lemma}[Feasibility]\label{lem:b1:greedy-proofs:feasible}
A set of tasks is feasible if and only if, for every $t$, at most $t$ of its tasks have a deadline at most $t$.
\end{lemma}

\begin{proof}
If more than $t$ tasks have a deadline at most $t$, they cannot all fit in the first $t$ slots. Conversely, run the
tasks by increasing deadline: the task in slot $i$ has the $i$-th smallest deadline, and at least $i$ tasks have a
deadline at most that one, so it is at least $i$.
\end{proof}

The greedy rule takes the tasks by decreasing penalty and keeps each one that leaves the set feasible.

\omcodes{code/bachelor-1/22-greedy-proofs/listings/penalties.py}{8-20}{The feasibility test and the greedy
choice}\label{lst:b1:greedy-proofs:penalties}

\omoutput{code/bachelor-1/22-greedy-proofs/outputs/penalties.txt}{Seven tasks, scheduled; and the greedy against all
subsets.}\label{out:b1:greedy-proofs:penalties}

\begin{theorem}[Greedy penalties]\label{thm:b1:greedy-proofs:penalties}
The greedy rule by decreasing penalty gives a feasible set of largest total penalty, so a schedule with the smallest
penalty paid.
\end{theorem}

\begin{proof}
First an exchange property: \emph{if $X$ and $Y$ are feasible and $|X| < |Y|$, some task $y$ of $Y$ not in $X$ keeps
$X + \{y\}$ feasible.} Write $N_t(S)$ for the number of tasks of $S$ with deadline at most $t$, and let $k$ be the
largest $t$ with $N_t(Y) \leq N_t(X)$ ($k = 0$ qualifies). For large $t$, $N_t(Y) = |Y| > |X|$, so $N_t(Y) > N_t(X)$
for every $t > k$, and $Y$ has more tasks than $X$ with a deadline after $k$: some $y \in Y - X$ has $d_y > k$. Adding
it to $X$ changes $N_t$ only for $t \geq d_y > k$, where $N_t(X) + 1 \leq N_t(Y) \leq t$: still feasible by the lemma.

Let $g_1, g_2, \dots$ be the greedy's tasks and $a_1, a_2, \dots$ those of an optimal set, both by decreasing penalty.
A subset of a feasible set is feasible, so whenever the greedy rejected a task, adding it to all the tasks it ends
with would not be feasible either. If the optimal set had more tasks, the exchange property would give a task that can be
added to the greedy's set, which the greedy would have kept: so $|A| \leq |G|$. Suppose $p(a_i) > p(g_i)$ for some $i$,
the first one. With $X = \{g_1, \dots, g_{i-1}\}$ and $Y = \{a_1, \dots, a_i\}$, some $a_j \notin X$, $j \leq i$, keeps
$X + \{a_j\}$ feasible, and $p(a_j) \geq p(a_i) > p(g_i)$. The greedy looked at $a_j$ before $g_i$, when its set was part
of $X$, so it would have kept it: then $a_j$ would be one of $g_1, \dots, g_{i-1}$, which it is not. So $p(g_i) \geq
p(a_i)$ for all $i \leq |A|$, and the greedy's total is at least the optimum.
\end{proof}

\section{Huffman codes}\label{sec:b1:greedy-proofs:huffman}

Text stored with a fixed number of bits per character wastes space on the frequent ones. The second book compressed
texts by giving short codes to frequent symbols; the question here is the best such code.

\begin{definition}[Prefix code, code tree, cost]\label{def:b1:greedy-proofs:prefix-code}
A \emph{prefix code}\index{prefix code} assigns to each symbol a word of bits, its \emph{codeword}, so that no codeword is
the beginning of another. Its \emph{code tree}\index{code tree} is the \omterm{def:b1:induction-and-termination:binary-tree}{binary tree} whose leaves are the symbols, the
codeword of a symbol spelling the path from the root, 0 for left and 1 for right. For frequencies $f_s$, the
\emph{cost}\index{cost (of a code)} of the code is $\sum_s f_s \ell_s$, where $\ell_s$, the length of the codeword of
$s$, is the depth of its \omterm{def:b1:induction-and-termination:rooted-tree}{leaf}: the number of bits of the encoded text.
\end{definition}

Without the prefix property, the bits would need separators: with $\mathtt{e} = 0$, $\mathtt{t} = 01$ and
$\mathtt{h} = 1$, the bits $01$ read either \code{t} or \code{eh}. With it, the decoder walks down the tree and outputs
a symbol at each \omterm{def:b1:induction-and-termination:rooted-tree}{leaf}.

\begin{method}[Huffman's algorithm]\label{met:b1:greedy-proofs:huffman}
Start with one tree per symbol, of weight its frequency. While there are two trees or more, remove the two lightest
and join them under a new root, of weight the sum of theirs. The last tree is the \omterm{def:b1:greedy-proofs:prefix-code}{code tree} of a \emph{\omterm{met:b1:greedy-proofs:huffman}{Huffman code}}.
\index{Huffman code}\index{Huffman's algorithm}
\end{method}

\omcodes{code/bachelor-1/22-greedy-proofs/listings/huffman.py}{18-37}{\omterm{met:b1:greedy-proofs:huffman}{Huffman's algorithm} with a \omterm{def:b1:heaps:pq}{priority queue}, and the
codewords}\label{lst:b1:greedy-proofs:huffman}

The \omterm{def:b1:heaps:pq}{priority queue} holds $(\mathit{weight}, \mathit{number}, \mathit{tree})$: the number breaks ties by age, so that two
runs on the same frequencies build the same tree, which matters to the decoder (\cref{exo:b1:greedy-proofs:6}). With $n$
symbols there are $n - 1$ merges, each $O(\log n)$.

\begin{omfigure}
\begin{tikzpicture}[n/.style={draw, circle, inner sep=1pt, minimum size=6mm, font=\scriptsize},
  l/.style={draw, rounded corners=2pt, inner sep=2pt, font=\scriptsize, fill=omDef!20},
  level distance=11mm, level 1/.style={sibling distance=48mm}, level 2/.style={sibling distance=24mm},
  level 3/.style={sibling distance=20mm}, level 4/.style={sibling distance=22mm},
  b/.style={font=\scriptsize, midway, fill=white, inner sep=0.5pt}]
\node[n] {43}
  child {node[n] {17}
    child {node[l] {a 8: 00} edge from parent node[b] {0}}
    child {node[l] {t 9: 01} edge from parent node[b] {1}} edge from parent node[b] {0}}
  child {node[n] {26}
    child {node[l] {e 12: 10} edge from parent node[b] {0}}
    child {node[n] {14}
      child {node[l] {o 7: 110} edge from parent node[b] {0}}
      child {node[n] {7}
        child {node[l] {s 3: 1110} edge from parent node[b] {0}}
        child {node[l] {n 4: 1111} edge from parent node[b] {1}} edge from parent node[b] {1}}
      edge from parent node[b] {1}} edge from parent node[b] {1}};
\node[font=\scriptsize, align=left, anchor=north west] at (-7.8,-0.2) {merges, in order:\\
1. s 3 + n 4 $\to$ 7\\ 2. o 7 + (7) $\to$ 14\\ 3. a 8 + t 9 $\to$ 17\\ 4. e 12 + (14) $\to$ 26\\
5. (17) + (26) $\to$ 43};
\end{tikzpicture}
\par\smallskip{\small \omterm{met:b1:greedy-proofs:huffman}{Huffman's algorithm} on six symbols of frequencies 12, 9, 8, 7, 4, 3: the merges, then the \omterm{def:b1:greedy-proofs:prefix-code}{code
tree}, each inner node marked with its weight, each \omterm{def:b1:induction-and-termination:rooted-tree}{leaf} with its symbol, frequency and codeword. The tie between o and
the tree of weight 7 goes to o, made first. \omterm{def:b1:greedy-proofs:prefix-code}{Cost}: $107$ bits, against $3 \times 43 = 129$ with three bits per symbol.}
\end{omfigure}

On the library's catalogue, the 80 titles in lower case, one per line:

\omoutput{code/bachelor-1/22-greedy-proofs/outputs/huffman.txt}{Frequencies and codewords of the catalogue's 24
symbols (\texttt{'\textbackslash n'} is the end of a line).}\label{out:b1:greedy-proofs:huffman}

\omcode{code/bachelor-1/22-greedy-proofs/listings/encode_decode.py}{12}{21}{Decoding by walking down the
tree}\label{lst:b1:greedy-proofs:decode}

\omoutput{code/bachelor-1/22-greedy-proofs/outputs/encode_decode.txt}{The first title, encoded and decoded; and the
whole catalogue.}\label{out:b1:greedy-proofs:decode}

\begin{lemma}[The two rarest are siblings]\label{lem:b1:greedy-proofs:siblings}
For any frequencies, some optimal \omterm{def:b1:greedy-proofs:prefix-code}{code tree} has the two least frequent symbols $x$ and $y$ as sibling leaves at the
largest depth.
\end{lemma}

\begin{proof}
In an optimal tree every inner node has two children: an inner node with one child could be replaced by that child,
shortening codewords. So a deepest \omterm{def:b1:induction-and-termination:rooted-tree}{leaf} has a sibling, also a \omterm{def:b1:induction-and-termination:rooted-tree}{leaf}; call them $a$ and $b$. Exchange $x$ with $a$: $x$ goes
down by $\ell_a - \ell_x \geq 0$ and $a$ goes up as much, and the \omterm{def:b1:greedy-proofs:prefix-code}{cost} changes by $(f_x - f_a)(\ell_a - \ell_x) \leq 0$,
since $f_x \leq f_a$. Exchange $y$ with $b$ in the same way. The tree is still optimal, and $x$, $y$ are sibling leaves
at the largest depth.
\end{proof}

\begin{lemma}[Merging keeps optimality]\label{lem:b1:greedy-proofs:merge}
Replace $x$ and $y$ by a new symbol $z$ of frequency $f_x + f_y$. If $T'$ is an optimal \omterm{def:b1:greedy-proofs:prefix-code}{code tree} for the new symbols,
the tree $T$ obtained by hanging $x$ and $y$ under the \omterm{def:b1:induction-and-termination:rooted-tree}{leaf} $z$ of $T'$ is optimal for the old ones.
\end{lemma}

\begin{proof}
The leaves $x$, $y$ of $T$ are one deeper than $z$ in $T'$, so $\mathrm{cost}(T) = \mathrm{cost}(T') + f_x + f_y$. Let
$U$ be an optimal tree for the old symbols with $x$, $y$ siblings (\cref{lem:b1:greedy-proofs:siblings}); replacing them
and their parent by a \omterm{def:b1:induction-and-termination:rooted-tree}{leaf} $z$ gives a tree $U'$ for the new symbols with $\mathrm{cost}(U) = \mathrm{cost}(U') + f_x +
f_y \geq \mathrm{cost}(T') + f_x + f_y = \mathrm{cost}(T)$. So $T$ is optimal.
\end{proof}

\begin{theorem}[Huffman]\label{thm:b1:greedy-proofs:huffman}
\omterm{met:b1:greedy-proofs:huffman}{Huffman's algorithm} builds a \omterm{def:b1:greedy-proofs:prefix-code}{prefix code} of smallest \omterm{def:b1:greedy-proofs:prefix-code}{cost}.
\end{theorem}

\begin{proof}
By induction on the number of symbols; one symbol needs a tree of one \omterm{def:b1:induction-and-termination:rooted-tree}{leaf}. The first merge joins $x$ and $y$, the two
lightest, and the rest of the algorithm is \omterm{met:b1:greedy-proofs:huffman}{Huffman's algorithm} on the symbols with $x$, $y$ replaced by $z$: by
induction it builds an optimal tree $T'$ for them, and \cref{lem:b1:greedy-proofs:merge} makes the final tree optimal.
\end{proof}

The proof never used the tie rule: any Huffman tree is optimal, though different tie rules give different codewords.
The construction was published by Huffman in 1952.

\section{When greedy is only an approximation}\label{sec:b1:greedy-proofs:approximation}

Change for an amount with the fewest coins: the greedy rule gives the largest coin that fits, repeatedly. With euro
coins it is optimal; with coins of 1, 3 and 4 it is not.

\omcode{code/bachelor-1/22-greedy-proofs/listings/coins.py}{4}{10}{Change, largest coin
first}\label{lst:b1:greedy-proofs:coins}

\omoutput{code/bachelor-1/22-greedy-proofs/outputs/coins.txt}{Greedy change against the fewest coins, for three coin
systems (amounts in cents).}\label{out:b1:greedy-proofs:coins}

The knapsack of \cref{sec:b1:backtracking:bb} can be filled greedily \omterm{def:b1:c-functions-and-stack:by-value}{by value} per unit of weight. That was the optimum of
the relaxed problem, where fractions are allowed; for whole items it is an approximation, good on average and bad on
some instances.

\omcode{code/bachelor-1/22-greedy-proofs/listings/knapsack_greedy.py}{6}{13}{The knapsack \omterm{def:b1:c-functions-and-stack:by-value}{by value} per unit of
weight}\label{lst:b1:greedy-proofs:knapsack}

\omoutput{code/bachelor-1/22-greedy-proofs/outputs/knapsack_greedy.txt}{The greedy knapsack against the optimum,
computed by the method of \cref{ch:b1:dynamic-programming}.}\label{out:b1:greedy-proofs:knapsack}

On the last instance, the greedy takes the small item of ratio 2 and has no room left for the item worth 100: its result
can be as far from the optimum as one likes. How far a greedy method can be from the optimum, in the worst case, is a
subject of the fifth book.

\section{Exercises}\label{sec:b1:greedy-proofs:exercises}

\begin{exercise}[$\star$]\label{exo:b1:greedy-proofs:1}
Which talks does the earliest-finish rule keep among $[2, 5)$, $[1, 3)$, $[4, 7)$, $[6, 8)$, $[3, 9)$, $[8, 10)$, $[7, 11)$,
$[9, 12)$? Can four talks fit?
\end{exercise}

\begin{exercise}[$\star$]\label{exo:b1:greedy-proofs:2}
Give the codeword lengths of a \omterm{met:b1:greedy-proofs:huffman}{Huffman code} for six symbols of frequencies 1, 1, 2, 3, 5, 8, and its \omterm{def:b1:greedy-proofs:prefix-code}{cost}.
\end{exercise}

\begin{exercise}[$\star$]\label{exo:b1:greedy-proofs:3}
Decode \code{111111001101110} with the \omterm{def:b1:greedy-proofs:prefix-code}{code tree} of the six-symbol figure.
\end{exercise}

\begin{exercise}[$\star$]\label{exo:b1:greedy-proofs:4}
Which coins does the greedy rule give for 289 cents with euro coins? For 10 with coins of 1, 3 and 4, and what is the
optimum?
\end{exercise}

\begin{exercise}[$\star\star$]\label{exo:b1:greedy-proofs:5}
This version sorts the talks by start. Explain its result, and say where the \omterm{met:b1:greedy-proofs:exchange}{exchange argument} breaks for it.
\omcode{code/bachelor-1/22-greedy-proofs/listings/exo_by_start.py}{4}{9}{Talks by start time}
\omoutput[2-3]{code/bachelor-1/22-greedy-proofs/outputs/exo_by_start.txt}{Five talks, two orders.}
\end{exercise}

\begin{exercise}[$\star\star$]\label{exo:b1:greedy-proofs:6}
Sender and receiver build the Huffman tree from the same frequencies, but break ties differently. Explain the output,
and give two ways to make decoding safe.
\omoutput[2-4]{code/bachelor-1/22-greedy-proofs/outputs/exo_ties.txt}{Four symbols of equal frequency.}
\end{exercise}

\begin{exercise}[$\star\star$]\label{exo:b1:greedy-proofs:7}
A station's timetable gives each train an arrival and a departure time; a train occupies a platform from its arrival to
its departure, both included. Compute the fewest platforms needed, in $O(n \log n)$.
\end{exercise}

\begin{exercise}[$\star\star$]\label{exo:b1:greedy-proofs:8}
Given only the codeword lengths of a \omterm{def:b1:greedy-proofs:prefix-code}{prefix code}, build the \emph{canonical} code: symbols sorted by length then by
name, each codeword the previous one plus one, followed by zeros when the length grows. Why is it enough for a sender to
transmit the lengths?
\end{exercise}

\begin{exercise}[$\star\star$]\label{exo:b1:greedy-proofs:9}
Show that the codeword lengths $\ell_1, \dots, \ell_n$ of a \omterm{def:b1:greedy-proofs:prefix-code}{prefix code} satisfy $\sum_i 2^{-\ell_i} \leq 1$, with equality
when every inner node of the \omterm{def:b1:greedy-proofs:prefix-code}{code tree} has two children.
\end{exercise}

\begin{exercise}[$\star\star\star$]\label{exo:b1:greedy-proofs:10}
Show that for the coins $1, c, c^2, \dots, c^k$, with $c \geq 2$, greedy change uses the fewest coins.
\end{exercise}

\begin{exercise}[$\star\star\star$]\label{exo:b1:greedy-proofs:11}
Tasks of lengths $t_i$ and deadlines $d_i$ must all be run, one after the other, on one machine; the lateness of a task
is its finishing time minus its deadline, if positive. Show that the order by earliest deadline minimises the largest
lateness, and test it against all orders.
\end{exercise}

\begin{exercise}[$\star\star\star$]\label{exo:b1:greedy-proofs:12}
Prove \cref{thm:b1:greedy-proofs:earliest-finish} again, with the other template: if $g_1, g_2, \dots$ are the greedy's
talks and $o_1, o_2, \dots$ those of an optimal set, both by finishing time, then $f(g_i) \leq f(o_i)$ for every $i$.
\end{exercise}

\section{Problem: Compressing the library catalogue}\label{sec:b1:greedy-proofs:problem}

\begin{problem}[{Weekend problem --- Compressing the library catalogue}]\label{pb:b1:greedy-proofs:1}
The catalogue is the text of \cref{out:b1:greedy-proofs:huffman}: the 80 titles in lower case, one per line. The problem
builds its \omterm{met:b1:greedy-proofs:huffman}{Huffman code} and proves that no \omterm{def:b1:greedy-proofs:prefix-code}{prefix code} does better.

\medskip\noindent\textbf{Part I --- A fixed-length code.}
\begin{enumerate}
\item How many characters and how many distinct symbols has the catalogue? Why is the end of line a symbol?
\item What is the shortest fixed length of codewords for its symbols?
\item How many bits does the fixed-length code take?
\item Which symbols are the most and the least frequent?
\item Is a fixed-length code a \omterm{def:b1:greedy-proofs:prefix-code}{prefix code}?
\item Show with an example why a code without the prefix property may not be decodable.
\end{enumerate}

\medskip\noindent\textbf{Part II --- Building the tree.}
\begin{enumerate}[resume]
\item Which two symbols does the first merge join, and what is the new weight?
\item The second merge has a choice between \code{g} and \code{u}, both of frequency 17. Which does the program take,
and why?
\item How many merges are made in all?
\item How many trees does the queue hold after $k$ merges?
\item What is the cost of the algorithm for $n$ symbols?
\item Why does the tie rule change the codewords but not the \omterm{def:b1:greedy-proofs:prefix-code}{cost}?
\end{enumerate}

\medskip\noindent\textbf{Part III --- Lengths, bits, decoding.}
\begin{enumerate}[resume]
\item Group the symbols by codeword length.
\item Check that $\sum_s 2^{-\ell_s} = 1$.
\item Compute the \omterm{def:b1:greedy-proofs:prefix-code}{cost} $\sum_s f_s \ell_s$ and compare with the program.
\item How many bits are saved against the fixed-length code, in percent?
\item Decode the first eleven bits of \cref{out:b1:greedy-proofs:decode} with the table.
\item Why does the decoder never need to read ahead?
\end{enumerate}

\medskip\noindent\textbf{Part IV --- Nothing does better.}
\begin{enumerate}[resume]
\item Which two symbols are siblings at the largest depth in the program's tree? Is this what
\cref{lem:b1:greedy-proofs:siblings} says?
\item Why does an optimal \omterm{def:b1:greedy-proofs:prefix-code}{code tree} have no inner node with one child?
\item In \cref{lem:b1:greedy-proofs:merge}, why is $\mathrm{cost}(T) = \mathrm{cost}(T') + f_x + f_y$?
\item Write the induction that proves \cref{thm:b1:greedy-proofs:huffman}.
\item Coding pairs of characters can take fewer bits than 5993. Why does this not contradict the theorem?
\item State the named result, \emph{Huffman's theorem on the catalogue}.
\end{enumerate}
\end{problem}
